\chapter{Phototherapy}

\section*{Overview}
Phototherapy is a treatment that exposes the skin to specific wavelengths of ultraviolet or visible light to treat skin disease, hyperbilirubinemia, or certain systemic disorders. It is delivered as narrowband or broadband UVB, UVA with photosensitizing agents (PUVA), or blue light for neonatal jaundice.

\section*{Alternative Names}
Light therapy; UV light therapy; ultraviolet phototherapy; PUVA; blue-light therapy; heliotherapy

\section*{Principle}
Light of specific wavelengths is absorbed by target molecules in the skin and produces therapeutic effects. In psoriasis and eczema, UVB and PUVA suppress epidermal hyperproliferation and modulate local immune responses. In neonatal jaundice, blue light within the 400 to 490 nm range converts unconjugated bilirubin to water-soluble photo-isomers that are excreted without conjugation.

\section*{History}
Phototherapy dates to ancient sunbathing traditions, and Niels Finsen won the Nobel Prize in 1903 for using light to treat tuberculosis of the skin. The use of blue light for neonatal jaundice was introduced in the 1950s.

\section*{Why It Is Done}
Ordered for moderate to severe psoriasis, vitiligo, atopic dermatitis, cutaneous T-cell lymphoma (mycosis fungoides), and other photosensitive skin disorders when topical therapy is inadequate. Blue-light phototherapy is ordered for significant neonatal unconjugated hyperbilirubinemia to prevent kernicterus.

\section*{Is This a Primary or Secondary Test or Procedure}
A primary treatment for neonatal hyperbilirubinemia and a second-line disease-modifying therapy for many dermatologic conditions.

\section*{How It Is Performed}
The patient stands in a cabinet lined with fluorescent tubes or LED arrays that emit the chosen wavelength while eyes are protected, with exposure times and doses titrated to skin type and response. PUVA requires taking a psoralen drug before exposure to sensitize the skin to UVA. In neonates, the infant is undressed, eyes are covered, and lights are positioned above the incubator or bassinet.

\section*{Preparation}
Skin type and minimal erythema dose are assessed, and photosensitizing drugs and natural sunlight are avoided where relevant. For PUVA, the psoralen is taken at a set time before exposure, and eye protection with UV-blocking glasses is used on treatment days. Bilirubin level is measured before starting infant phototherapy.

\section*{Contraindications and Precautions}
Conditions with photosensitivity, such as systemic lupus erythematosus, porphyria, and xeroderma pigmentosum, are contraindications, as is a history of melanoma or significant prior phototoxic exposure. Eye and gonadal protection is used, and careful dosimetry reduces the risk of burning; PUVA requires eye protection for the entire day of treatment.

\section*{Risks}
Short-term effects include erythema, pruritus, and photoaging of the skin. Long-term UVB and PUVA therapy increases the risk of skin cancer, and PUVA is associated with premature photoaging; infant phototherapy may cause loose stools, rash, and hyperthermia.

\section*{Aftercare and Recovery}
The dermatologic patient uses emollients, avoids sun exposure, and attends scheduled sessions, with doses adjusted to erythema response. In neonates, bilirubin levels are rechecked within hours and phototherapy is continued or intensified until levels fall below treatment thresholds, at which point it is stopped.

\section*{Outcome and Follow-up}
For skin disease, response is judged by reduction in lesions and symptom control. For neonatal jaundice, a falling bilirubin toward the age-specific treatment threshold indicates effective therapy.

\section*{Expected Results}
Clearance or significant improvement of skin lesions with dermatologic phototherapy, and bilirubin reduced to a safe level below the phototherapy threshold in the neonate.

\section*{Follow-up}
Periodic skin examinations are recommended for patients on long-term phototherapy because of skin cancer risk, and bilirubin monitoring continues until safe levels are sustained. Patients who do not respond may be offered phototherapy dose escalation, PUVA, or systemic therapy.

\section*{When to Seek Medical Care}
Follow the procedure team's specific instructions. Seek urgent medical assessment for severe or worsening pain, uncontrolled bleeding, fever or signs of infection, trouble breathing, chest pain, fainting, new weakness or confusion, inability to urinate, or any rapidly worsening symptom. The appropriate warning signs depend on the procedure and the patient's medical condition.

\section*{References}
\begin{enumerate}
  \item MedlinePlus. Phototherapy. U.S. National Library of Medicine; 2025.
  \item Current Medical Diagnosis and Treatment. McGraw-Hill; 2025.
\end{enumerate}

\clearpage
